\documentclass[11pt]{article}
\usepackage{amsmath,amssymb,amsthm,geometry}\geometry{margin=1in}
\newtheorem{theorem}{Theorem}\newtheorem{proposition}{Proposition}
\title{Progressive content substitution: grafting without a concatenation algebra}
\date{17 November 2026}\begin{document}\maketitle
\paragraph{Audit.} This construction includes native arrows, associativity and units. This remains a scoped construction, not a paid crossing. It consumes the progressive observation categories $\mathcal D_C(B)$ from \emph{Progressive Port Base and Unary Factorization Test}, and addresses the early-port kernel in \emph{Terminal Only Paid Factorization Obstruction}. Actual paid laws, coefficients and moving-input context actions remain independent obligations. Collections below are the original entire weighted ordered collections; no operation $C\circ C\to C$ is assumed.

\paragraph{Source route.} Read \texttt{\detokenize{Structured_Native_Extension_Family.tex}} for the paid-relative residual and \texttt{\detokenize{../globally-graded-open-interaction/Globally_Graded_Open_Interaction.tex}} for the original weighted-polynomial interface. The progressive base and content composition are companion files \texttt{\detokenize{Progressive_Port_Base_and_Unary_Factorization_Test.tex}} and \texttt{\detokenize{Progressive_Content_Substitution_Boundary.tex}} in this directory.

\section{Pointed content protocols and their dependent substitution}
Let $F:J\to\mathcal D_C(B)$ be a progressive content protocol. Its live ordered occurrences at $j$ are $P(j)$, with labels $b_p$. Its arrows preserve these labels by actual order-preserving injections. For every $b\in B$, supply a small child protocol $K_b$, a distinguished entry object $k_b^0$, and a functor $G_b:K_b\to\mathcal D_D(E)$, with one common entire collection $D$. Entries are typed data, not values sampled by content composition. Marked terminal objects are exactly those mapped to completed observations.

Define a dependent protocol $J[G]$. An object consists of $j\in J$ and one $k_p\in K_{b_p}$ for every live parent occurrence. An arrow to $(j',k')$ consists of an arrow $a:j\to j'$ and, writing $i:P(j)\hookrightarrow P(j')$, child arrows
\[
 k_p\to k'_{i(p)}\quad(p\in P(j)),\qquad
 k^0_{b_r}\to k'_r\quad(r\notin i(P(j))).                 \tag{1}
\]
For a composite, old-child arrows compose. If a child was born in the first arrow, its entry-to-child arrow composes with the second arrow's child arrow; if born only in the second arrow, retain that second entry arrow. This defines associative composition, since every occurrence has a unique first birth and subsequent child path. Identity parent arrows have no new occurrences and use identity child arrows. Equal parent endpoints do not identify different arrows or paid histories.

The entry-path field in (1) is material: a composite transition may contain both a birth and later child evolution. The elementary birth extension is the special case assigning a new child exactly $k_b^0$ with identity entry arrow. No stochastic or paid evolution is supplied by (1).

Define the output live ordered list by concatenating the $G_{b_p}(k_p)$ live lists in parent occurrence order. If the parent is not completed, or any child is not completed, mark the output unfinished. If the parent is completed with shape $c\in C_n$ and all children are completed with shapes $d_p\in D_{m_p}$, mark it completed with shape
\[
               (c;(d_p)_{p=1}^n)\in(C\circ D)_{\sum m_p}. \tag{2}
\]
The output list is the same concatenated occurrence list. In particular an unfinished child can expose outputs while its parent is still unfinished. A completed nullary parent gives a completed nullary output without children.

\begin{theorem}[Entire-grade progressive content substitution]
These formulas give a functor
\[
             F[G]:J[G]\longrightarrow\mathcal D_{C\circ D}(E).
\]
They retain all unused shapes and the actual coefficient objects
$C_c\otimes\bigotimes_pD_{d_p}$ of the entire substituted collection. Full arity-preserving collection maps act by their ordinary full substitution map and tensor of actual content arrows. No context-to-content arrow, paid coefficient, or concatenation algebra on $C$ is required.
\end{theorem}
\begin{proof}
Each old parent occurrence maps to its retained target occurrence and its child observation supplies an injection of old child positions. Their ordered disjoint union is an order-preserving injection of the flattened lists. New parent occurrences contribute only new output positions. Composition agrees positionwise with the composition in (1). If the source output is completed, the parent observation is completed and all child observations are completed. Every corresponding observation arrow is then the identity on its shape and list; no parent port can be born. Thus the target is completed with the same shape (2) and list, as required by $\mathcal D_{C\circ D}(E)$. Otherwise the constructed injection has a permitted unfinished source. This proves functoriality. The full grade is defined independently of which objects occur in $J[G]$, so unused shapes are not removed. Full collection-map and coefficient compatibility is precisely ordinary substitution functoriality, not an operation combining two shapes inside $C$.\end{proof}

\paragraph{Exact static restriction.}
Let $J$ be discrete with objects $S$, all completed, with $F(s)=(c_s,\mathbf b_s)$ the original static $N(S,B)$ body. For each $b$, let $K_b$ be the singleton entry protocol with completed value $G_b(k_b^0)=(d_b,\mathbf e_b)$, the original static $N(B,E)$ body. Then $J[G]$ is canonically discrete on $S$, and (2) is exactly original static content substitution, with its entire $C\circ D$ and actual content coefficients. Hence the progressive construction restricts to static $N$ without requiring arbitrary collections to be algebras for concatenation. This is an exact restriction, not an equivalence of every progressive protocol with a static body.

\section{Entry and moving-context boundaries}
\begin{proposition}[An unpointed child protocol does not determine birth]
There is no entry-selection operation natural for arbitrary unpointed child protocols, even when all content coefficients are units and all child outputs are terminal.
\end{proposition}
\begin{proof}
Take a discrete child protocol with two objects $u,v$ and a terminal content collection containing two nullary shapes, one assigned to each object. Simultaneously swapping the protocol objects and full grade shapes is an automorphism of this unpointed content datum. A natural choice of entry object would have to be fixed by this swap, which neither object is. Thus birth needs a supplied entry, an admitted family of entries or a weaker observer; a generic coend does not supply it.\end{proof}

For an actual random initializer, entry must instead be its full all-raw initialization interface, with actual law choices and once-owned coefficient factors. The pointed deterministic construction is not a theorem that those random entries commute with substitution. In particular replacing entry by an independently chosen random sample per occurrence can change whole-scope correlations and predecessor admission.

The reviewed open-response construction offers a relevant operation, but not an automatic solution: its states form a functorial passive action of a free observation-path category, and paid raw outcomes have subsequent input-indexed responses with all command/poststate squares. To use a moving $J$ here requires a context module indexed by $J$ (or an explicitly supplied functor from $J$ into that response base), together with transport of existing child states and the typed entry operation at every new occurrence. The transport must compose along $J$ arrows, preserve admission as declared, and commute with native initial/command/raw/coefficient/poststate fields. An incoming observation arrow is not by itself a paid event and must not introduce a new paid context factor. Arbitrary $J$ relations need not hold in the free response-path action, so assigning generators is insufficient without checking those relations.

Finally the category $J[G]$ retains all possible child observation paths, but supplies no causal scheduler or shared-world joint law to decide which such paths arise from a paid scope. Multiple child steps in (1) are paths, not a newly fused atomic draw. A paid crossing must separately construct the shape-blind context on this dependent protocol, prove exact owner/request whole-law admission, retain actual once-owned scope coefficients and full writer maps, and prove multiplication coherence. The content theorem resolves the entire-grade grafting and early-live-output part of that obligation only. It neither turns the representation by progressive bodies into a normal-form theorem nor identifies ordinary tensor and joint-limit paid contexts.

\section{Native arrows, reassociation and units}
Fix the intermediate label sets in this section. A native arrow of pointed protocols consists of an entry-preserving functor $u:K\to K'$ (strictly $u(k^0)=k'^0$), a full arity-preserving collection map $\beta:D\to D'$, and a natural transformation
\[
 \mathcal D_\beta(E)G\Longrightarrow G'u
\]
whose components are bijections on live occurrences and preserve labels. An unfinished source may become completed with the same list; at a completed source the full shape/body equation is required. All actual content arrows at unused shapes remain part of $\beta$. Use the corresponding definition without entry preservation for unpointed outer protocols. When a family of children is used, its members use the same full collection map $\beta$. This is a declared native arrow class, not all arbitrary functors of observation categories.

An outer native arrow and a family of child native arrows induce a functor on dependent protocols: map the outer object, reindex its child tuple by the actual parent-port bijection, and apply the corresponding child functors. On arrows use their supplied functor maps, including the entry paths in (1). Naturality of the outer port transformation identifies old and new occurrences; strict entry preservation types the images of newborn paths. The output transformation is the flattened ordered union of child port bijections. If its source output is completed, all source parent/child outputs are completed, so all terminal shape equations hold and give the full substitution map $\alpha\circ\beta:C\circ D\to C'\circ D'$. If its source is unfinished the same flattened bijection is an allowed progressive component. This proves substitution is functorial on these independently defined native arrows, including completion collapse. Componentwise composition proves bifunctoriality; no paid context map has been inferred.

Now supply a third family $H_e:L_e\to\mathcal D_V(Z)$ with entries $l_e^0$, where $V$ denotes the third entire weighted collection and $e$ ranges over the preceding output-label set. Each child composite $K_b[H]$ is pointed by
\[
        (k_b^0,(l^0_{e_r})_{r\in P(G_b(k_b^0))}).        \tag{3}
\]
Thus a child entry with already live ports does not lose their grandchildren. There is a canonical protocol isomorphism
\[
             (J[G])[H]\ \cong\ J[G[H]].                \tag{4}
\]
On objects it rebrackets exactly the same tuple $(j,(k_p)_p,(l_{p,r})_{p,r})$. On arrows it retains the outer arrow, each child arrow, and each grandchild arrow. To verify the nontrivial birth case, let $p$ be newly born along the outer arrow. The left side retains a path $k_b^0\to k'_p$, and, for every final child occurrence, an entry path $l_e^0\to l'_{p,r}$. On the right an arrow from (3) has the same child path. For child occurrences already present at $k_b^0$, its grandchild arrows start at the corresponding $l_e^0$; for occurrences newly born along the child path they also start at $l_e^0$. Hence the right side retains exactly the same arrows, not just their endpoints. For an old $p$, old and newly born grandchild occurrences are likewise identical on both sides. These cases exhaust all occurrences and prove the inverse bijection on arrows. Composition preserves first-birth paths on either side, so (4) is a category isomorphism.

The observation lists under (4) differ only by associativity of ordered concatenation. Completion on either side means exactly that the outer parent and every child and grandchild is completed. Its full terminal grade comparison is the ordinary entire-collection associator
\[
                   (C\circ D)\circ V\cong C\circ(D\circ V).
\]
For each shape, the coefficient arrow is the actual monoidal rebracketing (and the required symmetry for matching traversal orders) of the same $C$, $D$ and $V$ content factors. It is defined at every unused shape as well. No coefficient is duplicated or discarded. Native naturality follows on each retained tuple/path and on each full grade component.

For four layers every composite associator sends a nested configuration to the same flattened addressed tree of objects and paths. Each occurrence keeps its unique first-birth entry path; therefore the pentagon commutes on protocol objects and arrows, not merely output lists. On full weighted grades it is the ordinary substitution pentagon, and on actual content factors it is symmetric monoidal coherence. Consequently the native pentagon commutes at the stated interface.

The right unit family at a label $b$ is the singleton pointed protocol whose observation is the completed unary unit shape with sole label $b$. Substituting it retains exactly the original parent objects, arrows and live occurrences, giving $J[I]\cong J$ and the full right unitor $C\circ I\cong C$. For the left unit at $b$, take the singleton outer protocol with one completed unary occurrence labelled $b$. Its substitution by $K_b$ has precisely the objects and arrows of $K_b$, with entry $k_b^0$ when entries are retained, and full grade $I\circ D\cong D$. These are genuine unit isomorphisms. The triangle follows by the same addressed-tuple calculation and the ordinary full-grade unit coherence.

This supplies independent pointed progressive content multiplication, native associativity and units. Its protocol objects encode decoration and allowed observation paths only, not paid raw laws, world updates, costs or a scheduler. The construction includes all declared compatible path combinations, including paths not executable under a given paid context. Realizing such a protocol by fresh static processes still requires explicit installation at births and a separately typed moving-observation context action. Thus this content multiplication is useful structure for a paid relative candidate, not a proof of its crossing or of its normal forms.
\end{document}
